\oldpage[III]{190}
\subsection{The case of Noetherian local rings}
\label{subsection:III.7.5}

\begin{env}[7.5.1]
\label{III.7.5.1}
Let $A$ be a Noetherian local ring and $\mathfrak m$ its
maximal ideal.  For every $A$-module $M$, denote by
$\widehat M$ its separated completion for the
$\mathfrak m$-preadic topology, which is isomorphic to $\varprojlim\bigl(M\otimes_A(A/\mathfrak m^{n+1})\bigr)=\varprojlim\bigl(M/\mathfrak m^{n+1}M\bigr)$.
Let $T$ be a covariant additive functor from $\cat{Ab}_A$ to
$\cat{Ab}$.  The canonical homomorphisms
\sref{III.7.2.3.1}
\[
  T(M)\otimes_A(A/\mathfrak m^{n+1})
  \longrightarrow
  T\bigl(M\otimes_A(A/\mathfrak m^{n+1})\bigr)
\]
plainly form a projective system of $A$-homomorphisms.  On
passing to the limit they therefore give an
$\widehat A$-homomorphism, functorial in $M$,
\[
  (T(M))^{\widehat{\,}}
  \longrightarrow
  \varprojlim_n T(M_n),
  \tag{7.5.1.1}\label{III.7.5.1.1.eq}
\]
where $M_n=M\otimes_A(A/\mathfrak m^{n+1})$, $A_n=A/\mathfrak m^{n+1}$.
\end{env}

\begin{proposition}[7.5.2]
\label{III.7.5.2}
Let $A$ be a Noetherian local ring with maximal ideal
$\mathfrak m$ and residue field $k=A/\mathfrak m$.  Let $T$
be a covariant additive functor from $\cat{Ab}_A$ to
$\cat{Ab}$ that is semi-exact and commutes with inductive
limits.  Suppose, moreover, that for every finite-type
$A$-module $M$, the module $T(M)$ is of finite type and the
canonical homomorphism \eqref{III.7.5.1.1.eq} is an
isomorphism.  Then the following conditions are equivalent:
\begin{enumerate}
  \item[(a)] $T$ is right exact;
  \item[(b)] for every $n$, the functor $N\leadsto T(N)$ is
    right exact in the category of finite-type $A_n$-modules
    \textup{(which amounts to saying that $T$ is right exact in the
    category of $A$-modules \emph{of finite length})};
  \item[(c)] the canonical homomorphism $T(\varepsilon):T(A)\longrightarrow T(k)$
    is surjective;
  \item[(d)] for every sufficiently large $n$, the canonical
    homomorphism $T(A_n)\longrightarrow T(k)$
    is surjective.
\end{enumerate}
\end{proposition}

\begin{proof}
It is clear that \emph{(a)} implies \emph{(b)}.  Let us prove the converse.
Suppose that $u:M\to N$ is an epimorphism of $A$-modules; we
must show that $T(u)$ is an epimorphism.  Since $T$ commutes
with inductive limits and filtered inductive limits are exact
in the category of modules, it is enough to consider the case
where $M$ and $N$ are of finite type.

The modules $T(M)$ and $T(N)$ are then of finite type.  Since
$A$ is Noetherian and local, it is enough to prove that
\[
  (T(u))^{\widehat{\,}}:
  (T(M))^{\widehat{\,}}
  \longrightarrow
  (T(N))^{\widehat{\,}}
\]
is surjective \sref[0]{0.7.3.5} and \sref[0]{0.6.4.1}.
By hypothesis, $(T(M))^{\widehat{\,}}$ and $(T(N))^{\widehat{\,}}$ are,
respectively,
\[
  \varprojlim T(M_n)
  \quad\hbox{and}\quad
  \varprojlim T(N_n).
\]
Thus $(T(u))^{\widehat{\,}}$ is the inverse limit of the projective
system of homomorphisms $T(u\otimes1_{A_n}):T(M_n)\longrightarrow T(N_n)$.
Now \emph{(b)} means that these homomorphisms are
\emph{surjective}; moreover, $T(M_n)$ is a finite-type
$A_n$-module, and $A_n$ is an \emph{Artinian} ring by hypothesis.  It follows that
$(T(u))^{\widehat{\,}}$ is surjective \sref[0]{0.13.1.2} and
\sref[0]{0.13.2.2}.

It is clear that \emph{(a)} implies \emph{(c)}, and since $T(\varepsilon)$
factors as
$T(A)\longrightarrow T(A_n)\longrightarrow T(k)$,
\emph{(c)} implies \emph{(d)}; finally, it follows from \sref{III.7.2.7}
that \emph{(b)} and \emph{(d)} are equivalent, since $T$ is semi-exact in
$\cat{Ab}_{A_n}$, which completes the proof.
\end{proof}

\begin{corollary}[7.5.3]
\label{III.7.5.3}
Under the general hypotheses of \sref{III.7.5.2}, if
$T(k)=0$, then $T(M)=0$ for every $A$-module $M$.
\end{corollary}

\begin{proof}
\oldpage[III]{191}
Since $k$ is the only simple $A$-module,
\eqref{III.7.3.5.4.eq} shows that $T(E)=0$ for every
finite-length $A$-module $E$.  If $M$ is of finite type, then
\[
  (T(M))^{\widehat{\,}}
  \simeq
  \varprojlim_n T(M_n).
\]
Each $M_n$ has finite length, so $(T(M))^{\widehat{\,}}=0$.  By
hypothesis $T(M)$ is of finite type, and it embeds in
$(T(M))^{\widehat{\,}}$ \sref[0]{0.7.3.5}; hence $T(M)=0$.
Finally, for an arbitrary $A$-module $M$, the module $T(M)$
is the inductive limit of the $T(N_\alpha)$, where the
$N_\alpha$ run through the finite-type submodules of $M$.
The proof is complete.
\end{proof}

\begin{proposition}[7.5.4]
\label{III.7.5.4}
Let $A$ be a Noetherian local ring with maximal ideal
$\mathfrak m$ and residue field $k=A/\mathfrak m$.  Let
$T_\bullet$ be a homological functor from $\cat{Ab}_A$ to
$\cat{Ab}$ that commutes with inductive limits.  Suppose,
moreover, that for every $i$ and every finite-type $A$-module
$M$, the module $T_i(M)$ is of finite type and the canonical
homomorphism $(T_i(M))^{\widehat{\,}}\longrightarrow\varprojlim T_i(M_n)$
is bijective.  For a given integer $p$, the following
conditions are equivalent:
\begin{enumerate}
  \item[(a)] $T_p$ is exact;
  \item[(b)] $T_p$ is right exact and $T_p(A)$ is a free
    $A$-module;
  \item[(c)] the canonical homomorphisms
    \[
      T_{p+1}(A)\longrightarrow T_{p+1}(k),
      \qquad
      T_p(A)\longrightarrow T_p(k)
    \]
    are surjective;
  \item[(d)] for every $n$, the canonical homomorphisms
    \[
      T_{p+1}(A_n)\longrightarrow T_{p+1}(k),
      \qquad
      T_p(A_n)\longrightarrow T_p(k)
    \]
    are surjective;
  \item[(e)] for every $n$, the functor $N\leadsto T_p(N)$ is
    exact in the category of finite-type $A_n$-modules.
\end{enumerate}
\end{proposition}

\begin{proof}
By \sref{III.7.3.3}, condition \emph{(a)} is equivalent to saying
that $T_{p+1}$ and $T_p$ are right exact.  Since
$T_\bullet$ is a homological functor in $\cat{Ab}_{A_n}$, the
same argument as in \sref{III.7.3.1} shows that \emph{(e)} is
equivalent to saying that $T_p$ and $T_{p+1}$ are right exact
on the category of finite-type $A_n$-modules.  It follows from
\sref{III.7.5.2} that \emph{(a)} and \emph{(e)} are equivalent; the same
proposition also gives the equivalence of \emph{(a)}, \emph{(c)}, and \emph{(d)}.

Finally, every finite-type flat module over the local ring
$A$ is \emph{free} (Bourbaki, \emph{Alg\`ebre commutative},
chap.~II, \S~3, no.~2, cor.~2 of prop.~5).  The equivalence
of \emph{(a)} and \emph{(b)} therefore follows from \sref{III.7.3.1} and
\sref{III.7.3.3}.
\end{proof}

\begin{corollary}[7.5.5]
\label{III.7.5.5}
Suppose that the general hypotheses of \sref{III.7.5.4} hold.
\begin{enumerate}
  \item[\textup{(i)}] If $T_p(k)=0$, then $T_p=0$, the functor
    $T_{p+1}$ is right exact, and $T_{p-1}$ is left exact.
  \item[\textup{(ii)}] If $T_{p-1}(k)=T_{p+1}(k)=0$, then $T_p$ is
    exact, the canonical homomorphism
    \[
      T_p(A)\otimes_A M
      \longrightarrow
      T_p(M)
    \]
    is bijective, and $T_p(A)$ is a free $A$-module.
\end{enumerate}
\end{corollary}

\begin{proof}
Assertion (i) follows at once from \sref{III.7.5.3}, since
$T_p$ is semi-exact; the last assertion follows from the
definition of a homological functor.  Taking
\sref{III.7.3.3} into account, (i) immediately gives the first
two assertions of (ii), and the freeness of $T_p(A)$ follows
from \sref{III.7.5.4}.
\end{proof}

\begin{corollary}[7.5.6]
\label{III.7.5.6}
Suppose that the general hypotheses of \sref{III.7.5.4} hold,
and suppose in addition that $A$ is a discrete valuation
ring.
\begin{enumerate}
  \item[\textup{(i)}] The functor $T_p$ is right exact if and only if
    $T_{p-1}(A)$ is a free $A$-module.
  \item[\textup{(ii)}] The functor $T_p$ is exact if and only if
    $T_p(A)$ and $T_{p-1}(A)$ are free $A$-modules.
\end{enumerate}
\end{corollary}

\oldpage[III]{192}
\begin{proof}
It is clear that (ii) implies (i) (cf.\
\sref{III.7.3.3}).  To prove (i), let $f$ be a generator of
the maximal ideal of $A$, that is, a uniformizer of $A$.  A
finite-type $A$-module $M$ is free, equivalently flat, if and
only if the homothety
\[
  h_f:M\longrightarrow M,\qquad x\to fx,
\]
is injective, since in the present situation this amounts to
saying that $M$ is torsion-free \sref[0]{0.6.3.4}.  The exact
sequence $0\longrightarrow A\xrightarrow{h_f}A\longrightarrow k\longrightarrow0$
gives the exact homology sequence
\[
  T_p(A)
  \longrightarrow T_p(k)
  \longrightarrow T_{p-1}(A)
  \xrightarrow{h_f}
  T_{p-1}(A)
\]
Thus $T_{p-1}(A)$ is free if and only if
$T_p(A)\to T_p(k)$ is surjective, and the conclusion follows
from \sref{III.7.5.2}.
\end{proof}

\begin{remark}[7.5.7]
\label{III.7.5.7}
By \sref{III.7.4.7}, the general hypotheses of
\sref{III.7.4.4} imply that the homological functor
$T_\bullet$ defined by \eqref{III.7.4.1.1.eq} satisfies the
general hypotheses of \sref{III.7.5.4}.  In that case,
\sref{III.7.5.6} is therefore contained in
\sref{III.7.4.3}.
\end{remark}
